% reality/realized_pose3.tex

\chapter{Realized Pose3 and Dynamic Realization}

% ============================================================
\section{Motivation}
% ============================================================

\begin{remark}
Classical alignment theory typically assumes that the designed geometry
is identical to the realized geometry.
\end{remark}

\begin{remark}
In reality, however, railway tracks are produced through iterative and
physically constrained realization processes.
\end{remark}

\begin{remark}
As a consequence, not only the geometric position of the track differs
from the theoretical target alignment, but also its orientation, cant,
and dynamic behaviour.
\end{remark}

\begin{remark}
Thus, realization affects the full spatial railway state rather than
geometry alone.
\end{remark}

\begin{remark}
This motivates the introduction of a realized pose3 state:
\[
\mathrm{pose3}_{\mathrm{real}}(s).
\]
\end{remark}

% ============================================================
\section{Target and Realized States}
% ============================================================

\begin{definition}[Target pose3]
The target railway state is defined as
\[
\mathrm{pose3}_{\mathrm{target}}(s)
=
\bigl(
\gamma_{\mathrm{target}}(s),
\theta_{\mathrm{target}}(s),
\phi_{\mathrm{target}}(s)
\bigr).
\]
\end{definition}

\begin{definition}[Realized pose3]
The realized railway state is defined as
\[
\mathrm{pose3}_{\mathrm{real}}(s)
=
\bigl(
\gamma_{\mathrm{real}}(s),
\theta_{\mathrm{real}}(s),
\phi_{\mathrm{real}}(s)
\bigr).
\]
\end{definition}

\begin{remark}
The realized state is the physically constructed and measurable state of
the railway track.
\end{remark}

\begin{remark}
In practice,
\[
\mathrm{pose3}_{\mathrm{real}}
\neq
\mathrm{pose3}_{\mathrm{target}}.
\]
\end{remark}

% ============================================================
\section{Sources of Realization Deviations}
% ============================================================

\begin{remark}
Deviations between target and realized state arise from:
\begin{itemize}
\item measurement uncertainty,
\item machine limitations,
\item discrete control points,
\item structural elasticity,
\item ballast behaviour,
\item rail bending,
\item settlement processes,
\item and operational degradation.
\end{itemize}
\end{remark}

\begin{remark}
Thus, realization is not a purely geometric transformation, but a
mechanical and operational process.
\end{remark}

% ============================================================
\section{Realized Curvature and Cant}
% ============================================================

\begin{definition}[Realized curvature]
The realized curvature is
\[
\kappa_{\mathrm{real}}(s)
=
\frac{\dd \theta_{\mathrm{real}}}{\dd s}.
\]
\end{definition}

\begin{definition}[Realized cant]
The realized cant angle is
\[
\phi_{\mathrm{real}}(s).
\]
\end{definition}

\begin{remark}
Both quantities differ from their target values due to realization
effects.
\end{remark}

\begin{remark}
In particular,
\[
\kappa_{\mathrm{real}}
\neq
\kappa_{\mathrm{target}},
\qquad
\phi_{\mathrm{real}}
\neq
\phi_{\mathrm{target}}.
\]
\end{remark}

% ============================================================
\section{Cant Realization}
% ============================================================

\begin{remark}
Cant realization is particularly sensitive because cant is not only a
geometric quantity but also a dynamically relevant operational state.
\end{remark}

\begin{remark}
Cant realization is affected by:
\begin{itemize}
\item tamping accuracy,
\item rail fastening systems,
\item ballast stiffness,
\item torsion limits,
\item local support conditions,
\item and measurement procedures.
\end{itemize}
\end{remark}

\begin{remark}
In practice, cant ramps are realized only approximately and may exhibit
local irregularities.
\end{remark}

\begin{remark}
This is especially relevant in:
\begin{itemize}
\item transition zones,
\item switches and crossings,
\item station areas,
\item and overlapping transition regions.
\end{itemize}
\end{remark}

% ============================================================
\section{Coupling of Curvature and Cant}
% ============================================================

\begin{remark}
Curvature and cant are dynamically coupled quantities.
\end{remark}

\begin{remark}
Therefore, realization errors in curvature and cant cannot be treated
independently.
\end{remark}

\begin{remark}
The physically relevant quantity is the realized dynamic state:
\[
(\kappa_{\mathrm{real}}, \phi_{\mathrm{real}}).
\]
\end{remark}

\begin{definition}[Realized equilibrium condition]
A realized track is dynamically balanced if
\[
v^2\kappa_{\mathrm{real}}(s)
=
g\sin\phi_{\mathrm{real}}(s).
\]
\end{definition}

\begin{remark}
In practice, deviations from equilibrium are unavoidable and constrained
only within admissible limits.
\end{remark}

% ============================================================
\section{Dynamic Realization Error}
% ============================================================

\begin{definition}[Realized effective lateral acceleration]
The realized effective lateral acceleration is
\[
a_{\mathrm{eff,real}}(s)
=
v^2\kappa_{\mathrm{real}}(s)
-
g\sin\phi_{\mathrm{real}}(s).
\]
\end{definition}

\begin{remark}
This quantity is operationally more relevant than geometric deviation
alone.
\end{remark}

\begin{definition}[Realized jerk]
The realized jerk is
\[
j_{\mathrm{real}}(s)
=
v^3\kappa'_{\mathrm{real}}(s)
-
gv\cos\phi_{\mathrm{real}}(s)\,
\phi'_{\mathrm{real}}(s).
\]
\end{definition}

\begin{remark}
Passenger comfort depends directly on the realized jerk rather than on
the analytical target functions.
\end{remark}

% ============================================================
\section{Realization as State Transformation}
% ============================================================

\begin{definition}[Pose3 realization operator]
The realization process may be written as
\[
\mathcal{R}_{\mathrm{pose3}}
:
\mathrm{pose3}_{\mathrm{target}}
\mapsto
\mathrm{pose3}_{\mathrm{real}}.
\]
\end{definition}

\begin{remark}
This operator includes:
\begin{itemize}
\item geometric realization,
\item cant realization,
\item mechanical filtering,
\item and operational constraints.
\end{itemize}
\end{remark}

\begin{remark}
Thus, realization acts on the full railway state rather than only on
position space.
\end{remark}

% ============================================================
\section{Filtering Interpretation}
% ============================================================

\begin{proposition}
The realization operator acts as a low-pass filter on:
\begin{itemize}
\item curvature,
\item cant,
\item and higher-order derivatives.
\end{itemize}
\end{proposition}

\begin{remark}
High-frequency components are attenuated through:
\begin{itemize}
\item machine limitations,
\item structural smoothing,
\item ballast mechanics,
\item and measurement discretization.
\end{itemize}
\end{remark}

\begin{remark}
Therefore, highly detailed analytical target states may collapse into
similar realized states.
\end{remark}

% ============================================================
\section{Operational Interpretation}
% ============================================================

\begin{remark}
Railway operation interacts with the realized state rather than with the
theoretical target geometry.
\end{remark}

\begin{remark}
Consequently, the relevant engineering question is not:
\begin{quote}
``Was the theoretical geometry built exactly?''
\end{quote}
but rather:
\begin{quote}
``Does the realized state produce the intended dynamic behaviour?''
\end{quote}
\end{remark}

% ============================================================
\section{Connection to Optimization}
% ============================================================

\begin{remark}
Optimization should therefore operate on realized states:
\[
J
=
J\bigl(
\mathrm{pose3}_{\mathrm{real}}
\bigr)
=
J\bigl(
\mathcal{R}_{\mathrm{pose3}}
(
\mathrm{pose3}_{\mathrm{target}}
)
\bigr).
\]
\end{remark}

\begin{remark}
This extends classical alignment optimization toward realization-aware
optimization.
\end{remark}

% ============================================================
\section{Connection to AIM}
% ============================================================

\begin{summary}
Real railway tracks are not analytical curves but realized physical
states.

The realization process acts on the complete pose3 state, including
geometry, curvature, and cant.

Thus, railway alignment design must be understood as the design of a
realizable dynamic state rather than merely a geometric target curve.

Within the AIM framework, realized pose3 forms the bridge between
geometry, mechanics, construction, and operation.
\end{summary}
